\section*{الفصل \ref{ch:b2:angiosperm-reproduction} --- التكاثر الجنسي عند النباتات الزهرية}

\begin{solution}{exo:b2:angiosperm-reproduction:1}
السبلات (الكأس)، والبتلات (التويج)، والسدى (الطقم الذكري)، والكرابل
(الطقم الأنثوي). \omterm{def:b2:angiosperm-reproduction:flower}{والسداة}: خيط ومتك. \omterm{def:b2:angiosperm-reproduction:flower}{والكربلة}: \omterm{def:b2:angiosperm-reproduction:flower}{ميسم} وقلم ومبيض فيه
بييضات. وكل هذه \omterm{def:b2:plant-life-cycles:cycles}{طور بوغي} \omterm{def:b2:meiosis-heredity:meiosis}{ثنائي الصيغة}؛ ولا يكون طورًا مشيجيًا فرديَ
الصيغة إلا \omterm{prop:b2:angiosperm-reproduction:pollen}{حبة اللقاح} (داخل المتك) \omterm{prop:b2:angiosperm-reproduction:embryosac}{والكيس الجنيني} (داخل \omterm{def:b2:plant-life-cycles:heterospory}{البييضة}).
\end{solution}

\begin{solution}{exo:b2:angiosperm-reproduction:2}
\omterm{def:b2:plant-life-cycles:heterospory}{البوغ الصغير} $n$؛ والخلية الخضرية $n$؛ والنطفة $n$؛ والبويضة $n$؛
والنواة القطبية $n$؛ واللاقحة $2n$؛ والسويداء $3n$؛ \omterm{def:b2:angiosperm-reproduction:seed}{وقشرة البذرة}
$2n$ (وهي الغلافان الأموميان)؛ وجدار \omterm{def:b2:angiosperm-reproduction:seed}{الثمرة} $2n$ (وهو المبيض الأمومي).
\end{solution}

\begin{solution}{exo:b2:angiosperm-reproduction:3}
الرياح (النجيليات، والبلوط، والبتولا، والبندق، ولسان الحمل): أزهار
خضراء صغيرة، بلا بتلات ولا رائحة ولا \omterm{def:b2:angiosperm-reproduction:pollination}{رحيق}، ومتوك متدلية على خيوط
طويلة، ومياسم ريشية، وكميات هائلة من \omterm{def:b2:plant-life-cycles:heterospory}{اللقاح} الأملس الجاف، وإزهار قبل
الأوراق. والنحل (النفل، والمريمية، وقفاز الثعلب، والخزامى، والتفاح):
بتلات ملونة عليها أدلة فوق بنفسجية غالبًا، ومنصة هبوط، ورائحة،
\omterm{def:b2:angiosperm-reproduction:pollination}{ورحيق}، \omterm{def:b2:plant-life-cycles:heterospory}{ولقاح} لزج منحوت بكميات معتدلة، وإزهار في ضوء النهار.
\end{solution}

\begin{solution}{exo:b2:angiosperm-reproduction:4}
تندمج خليتا النطفة اللتان \omterm{thm:b2:angiosperm-reproduction:tube}{لأنبوب لقاح} واحد، إحداهما بالبويضة
(اللاقحة، $2n$، ومنها الجنين) والأخرى بالخلية المركزية (نواة
السويداء، $3n$، ومنها مخزون الغذاء). ولا يبدأ السويداء بالنمو إلا
بعد الإلقاح، فلا يموّن النبات إلا بييضات تحمل جنينًا، بخلاف عاري
البذور الذي يبني مخزون غذائه قبل ذلك.
\end{solution}

\begin{solution}{exo:b2:angiosperm-reproduction:5}
$25/1.2 = \qty{21}{h}$. والحرائر البارزة في الأيام 5 و6 و7 تلقى لقاحًا؛
أما حرائر الأيام 8 إلى 12 فلا: أي إن $3/8$ من الحرائر تُلقّح، ويكون
الكوز خاليًا في خمسة أثمانه.
\end{solution}

\begin{solution}{exo:b2:angiosperm-reproduction:6}
$S_1S_2\times S_1S_2$: صفر. و$S_1S_2\times S_1S_3$: $\tfrac12$ (فلا ينمو إلا \omterm{def:b2:plant-life-cycles:heterospory}{لقاح}
$S_3$). و$S_1S_2\times S_3S_4$: الكل. وخلف التهجين
الثاني: البويضات $S_1$ أو $S_2$، والنطاف $S_3$: أي $S_1S_3$ و$S_2S_3$،
بالنصف لكل منهما.
\end{solution}

\begin{solution}{exo:b2:angiosperm-reproduction:7}
يحمل النبات 2 من الأليلات $n$ ويرفض \omterm{def:b2:plant-life-cycles:heterospory}{اللقاح} الحامل أيًا منهما:
فالكسر المتوافق $(n - 2)/n$. \omterm{def:b2:meiosis-heredity:vocabulary}{والأليل} الجديد لا ترفضه أي مدقة،
فيُنجب لقاحه بذورًا أكثر من المتوسط وينتشر حتى يصير شائعًا كالبقية؛
والميزة نفسها تحمي كل \omterm{def:b2:meiosis-heredity:vocabulary}{أليل} نادر من الضياع. فالانتخاب إذن يراكم
الأليلات، وتحمل الجمهرات الطبيعية للأنواع غير المتوافقة ذاتيًا عشراتها.
\end{solution}

\begin{solution}{exo:b2:angiosperm-reproduction:8}
\omterm{prop:b2:angiosperm-reproduction:embryosac}{النواتان القطبيتان} نسختان من \omterm{def:b2:plant-life-cycles:heterospory}{بوغ كبير} واحد، فتكون الخلية المركزية
$YY$ أو $yy$ (بالنصف لكل منهما)؛ والنطفة $Y$ أو $y$ (بالنصف لكل
منهما): فيكون السويداء $YYY$ أو $YYy$ أو $Yyy$ أو $yyy$ بنسبة $\tfrac14$ لكل منها؛ أي إن ثلاثة
أرباعها صفراء. وإذا لُقّح \omterm{def:b2:plant-life-cycles:heterospory}{بلقاح} $yy$: فالسويداء $YYy$ أو $yyy$، بالنصف لكل منهما ---
أي نصف الحبات صفراء.
\end{solution}

\begin{solution}{exo:b2:angiosperm-reproduction:9}
تُشق حبات الشعير نصفين؛ فتهضم الأنصاف الحاملة للجنين نشاءها، ولا
تهضمه الأنصاف الخالية منه، إلا إذا أُضيف الجبريلين فتهضمه. (أ)
فالجنين مصدر الإشارة لا مصدر الإنزيم. (ب) والجبريلين هو الإشارة،
وهو كافٍ بذاته. (ج) والأميلاز تصنعه طبقة الألورون من السويداء،
وهي التي تستجيب للهرمون. ونصف الحبة الخالي من الجنين يفصل الإشارة
عن الاستجابة: فبدونه لا يمكن معرفة هل يهضم الجنين النشاء بنفسه.
\end{solution}

\begin{solution}{exo:b2:angiosperm-reproduction:10}
الرياح: $10^{6}\times\qty{1}{\micro g} = \qty{1}{g}$ لكل \omterm{def:b2:plant-life-cycles:heterospory}{بييضة}. والحيوان:
\qty{1}{mg} من \omterm{def:b2:plant-life-cycles:heterospory}{اللقاح} زائد \qty{0.5}{mg} من السكر $= \qty{1.5}{mg}$
لكل \omterm{def:b2:plant-life-cycles:heterospory}{بييضة} --- أي أرخص بنحو سبعمئة مرة. لكن الإلقاح بالرياح لا
يحتاج إلى شريك: فهو يعمل في مطلع الربيع قبل أن تطير الحشرات، وفي
الأماكن الباردة والعاصفة والمكشوفة، وليلًا وفي المطر، وفي غابات
كثيفة من نوع واحد، ولا يستطيع سارقو \omterm{def:b2:angiosperm-reproduction:pollination}{الرحيق} خداعه ولا يهجره ملقّح
ينقرض.
\end{solution}

\begin{solution}{exo:b2:angiosperm-reproduction:11}
البذور التي تسقط عند قدم الأم تحط في ظلها، وبين جذورها، وحيث تتركز
مفترساتها البذرية وممرضاتها النوعية؛ فتموت كلها تقريبًا. أما الطائر
فيحمل \omterm{def:b2:angiosperm-reproduction:seed}{بذرة} كيلومترات
(\qty{12}{km} في عشرين دقيقة) ويسقطها، مع سماد، في
سياج أو فرجة؛ وحتى لو حطت \omterm{def:b2:angiosperm-reproduction:seed}{بذرة} واحدة من كل مئة في موضع جيد، فذلك
أكثر من كل البذور التي تحت الشجرة، وتنتشر الجمهرة بسرعة الطيور لا
بسرعة \omterm{def:b2:angiosperm-reproduction:seed}{الثمرة} الساقطة. فاللحم ثمن التذكرة.
\end{solution}

\begin{solution}{exo:b2:angiosperm-reproduction:12}
الإحاطة تجعل \omterm{def:b2:plant-life-cycles:heterospory}{اللقاح} ينمو عبر نسيج أمومي، وهذا يتيح للمدقة أن
تختبره (\omterm{prop:b2:angiosperm-reproduction:selfing}{عدم التوافق الذاتي}، ورفض الأنواع الغريبة)، وأن تدع أنابيب
كثيرة تتنافس فيُنجب أسرعها \omterm{def:b2:angiosperm-reproduction:seed}{البذرة}، وأن تحمي البييضات من الجفاف ومن
العواشب، وأن تحوّل المبيض ثمرةً تجنّد الحيوانات للانتشار؛ ثم يموّن
\omterm{prop:b2:angiosperm-reproduction:double}{الإلقاح المضاعف} البييضات الملقحة وحدها. أما الكلفة: فعلى \omterm{prop:b2:angiosperm-reproduction:pollen}{حبة اللقاح}
أن تبلغ ميسمًا لا \omterm{def:b2:plant-life-cycles:heterospory}{البييضة} نفسها، فيلزم قلم وأنبوب وهداية له، \omterm{def:b2:angiosperm-reproduction:seed}{والثمرة}
استثمار ثقيل. وقد رجح الميزان لكاسيات البذور رجحانًا شديدًا حتى
انتقلت من العدم إلى تسعة أعشار أنواع النباتات في مئة
مليون سنة.
\end{solution}

\begin{solution}{pb:b2:angiosperm-reproduction:1}
\textbf{1.} على A ($S_1S_2$): \omterm{def:b2:plant-life-cycles:heterospory}{لقاح} A يعطي $f = 0$؛ \omterm{def:b2:plant-life-cycles:heterospory}{ولقاح} B ($S_1$، $S_3$) يعطي $f =
\tfrac12$؛ \omterm{def:b2:plant-life-cycles:heterospory}{ولقاح} C ($S_3$، $S_4$) يعطي $f = 1$.
\textbf{2.} على B ($S_1S_3$): A يعطي $\tfrac12$، وB صفرًا، وC $\tfrac12$. وعلى C
($S_3S_4$): A يعطي 1، وB $\tfrac12$، وC صفرًا.
\textbf{3.} 50 حبة متوافقة؛ ومنها $0.1\times 50 = 5$ بييضات.
\textbf{4.} من C: 100 متوافقة، ومنها $0.1\times 100 = 10$، أي
البييضات العشر كلها؛ ومن A: لا شيء.
\textbf{5.} تُبقى الثمار الآتية من زيارات C (10 بذور)؛ والآتية من B (5 بذور)
تُبقى بالكاد؛ والآتية من A تسقط. وبستان من A وحده لا يثمر شيئًا.
\textbf{6.} نصف الأزهار يعقد ثمرًا (وهي المزارة من C): أي 2500
\omterm{def:b2:angiosperm-reproduction:seed}{ثمرة} في الشجرة، عشرة أضعاف الثمار 250 اللازمة --- فالمحصول كامل،
وستسقط الشجرة الفائض.
\textbf{7.} لأنها تنثر لقاحًا متوافقًا على مدى موسم إزهار كل
صنف؛ لكنها إن شاركت صنفًا في أليلي $S$ معًا لم تلقّح عليه شيئًا.
\textbf{8.} $8\times 10^{7}/(7\times 86400) = 132$ حبة في المتر
المربع في الثانية.
\textbf{9.} $132/0.25 = 530$ حبة في المتر المكعب.
\textbf{10.} $132\times 10^{-4} = 0.013$ في الثانية: أي حبة كل
\qty{76}{s}.
\textbf{11.} $0.013\times 86400 \approx 1100$ حبة في اليوم؛ وأول
أنبوب يبلغ \omterm{def:b2:plant-life-cycles:heterospory}{البييضة} يلقّحها، فتُغلق \omterm{def:b2:plant-life-cycles:heterospory}{البييضة} عندئذ أمام البقية.
\textbf{12.} بعد \qty{20}{h} من \omterm{def:b2:angiosperm-reproduction:pollination}{الإلقاح اللقاحي}.
\textbf{13.} \omterm{def:b2:plant-life-cycles:heterospory}{اللقاح} في الأيام 0--7، والحرائر من اليوم 5: فلا تتداخل
إلا الأيام 5 و6 و7 --- أي $3/8$ من الحرائر، والبارزة بعد
اليوم 7 لا تنال شيئًا؛ فيكون الكوز خاليًا في جله. والجفاف عند
الإحرار هو السبب الكلاسيكي لفشل محصول الذرة.
\textbf{14.} الجنين $Yy$؛ والسويداء $YYy$.
\textbf{15.} صفراء، لأن $Y$ \omterm{def:b2:meiosis-heredity:vocabulary}{سائد}: فصفة أبي \omterm{def:b2:plant-life-cycles:heterospory}{اللقاح} تظهر في \omterm{def:b2:angiosperm-reproduction:seed}{البذرة}
على النبات الأم --- وهي الغريبية.
\textbf{16.} الخلية المركزية $YY$ أو $yy$: فالسويداء $YYy$ (بالنصف) و
$yyy$ (بالنصف).
\textbf{17.} $YYy$: 20 وحدة؛ و$Yyy$: 10 وحدات؛ أما $yyY$ الذي جاء $Y$
فيه من الأب فهو \omterm{def:b2:meiosis-heredity:vocabulary}{النمط الوراثي} $Yyy$ نفسه، أي 10 وحدات. ولا يمكن أن
توجد جرعتان أبويتان: فنطفة واحدة تلقّح الخلية المركزية، والسويداء
دائمًا مورثومان أموميان إلى واحد أبوي.
\textbf{18.} لأنه مخزن النشاء؛ وهو $0.3/0.33 = \qty{91}{\%}$ من
الحبة.
\textbf{19.} لأن المخزون لا يُبنى إلا بعد الإلقاح، فلا تذهب موارد
إلى بييضات لا تحمل جنينًا؛ أما طور عاري البذور المشيجي فيُبنى قبل
ذلك ويُهدر إن فشل \omterm{def:b2:angiosperm-reproduction:pollination}{الإلقاح اللقاحي}.
\textbf{20.} $0.9\times 600 = 540$ حبة في الكوز وفي النبات؛ و
$4320$ في المتر المربع؛ و$4320\times 0.3 = \qty{1.3}{kg}$ من السويداء
في المتر المربع.
\textbf{21.} أي \qty{13}{t/ha}.
\textbf{22.} كربون الحبوب $0.5\times 1.2 = \qty{0.6}{kg/m^2}$؛
وكربون السويداء $0.45\times 1.3 = \qty{0.58}{kg/m^2}$: أي متسق
(والباقي الصغير جنين وقشرة).
\textbf{23.} $8\times 10^{7}/4320 \approx \num{18500}$ \omterm{prop:b2:angiosperm-reproduction:pollen}{حبة لقاح} لكل
حبة.
\textbf{24.} لأن سحابة لقاحه رقيقة وتذروها الرياح، ولأن نورته
الذكرية تنثر في معظمها قبل بروز حرائره، فلا تُلقّح حرائر كثيرة
ويكون في الكوز فجوات.
\textbf{25.} كتلة A: نصف الأزهار يعقد ثمرًا، أي 2500 في الشجرة، وهو
محصول كامل. والحقل: 4320 حبة في المتر المربع، أي \qty{13}{t/ha}.
\end{solution}
